\section{Interpolación polinómica}
\subsection{Motivación}

Hay un problema clásico en análisis numérico que es el siguiente:\\
\textit{Dado un conjunto de puntos $x_0, x_1, \ldots, x_n$ y un conjunto de valores $y_0, y_1, \ldots, y_n$, encontrar una función $f$ tal que $f(x_i) = y_i$ para todo $i = 0, 1, \ldots, n$.}
Para ello, una de las formas más comunes de resolver este problema es mediante la interpolación polinómica, que consiste en encontrar un polinomio $P$ de grado a lo sumo $n$ tal que $P(x_i) = y_i$ para todo $i = 0, 1, \ldots, n$.\\

\begin{definición}[Espacio de polinomios]
    El espacio de polinomios al conjunto:
    \[
    \mathbb{R}[x] = \{ p = a_0 + a_1 x + a_2 x^2 + \ldots + a_n x^n : a_i \in \mathbb{R}, n \in \mathbb{N} \}
    \]
    En particular, el espacio de polinomios de grado a lo sumo $n$ se denota por:
    \[
    \mathcal{P}_n = \{ p = a_0 + a_1 x + a_2 x^2 + \ldots + a_n x^n : a_i \in \mathbb{R} \}
    \]
\end{definición}

Un caso particular de este problema es el siguiente:

\[
    \mathcal{P}_0 = \{ a_0 : a_0 \in \mathbb{R} \}
\]

Que es isomorfo a $\mathbb{R}$, es decir:
\[
    \mathcal{P}_0 \cong \mathbb{R}
\]

De hecho, el espacio de polinomios de grado a lo sumo $n$ es isomorfo a $\mathbb{R}^{n+1}$, es decir:
\[
    \mathcal{P}_n \cong \mathbb{R}^{n+1}
\]

\subsection{Problema de la interpolación polinómica}

Formalizemos ahora el problema de la interpolación polinómica.

\begin{definición}[Problema de la interpolación polinómica]
    Sean:
    \[
        \begin{cases}
            \{ x_i \}_{i=0}^n \subset [a,b] : \quad x_i \neq x_j \iff i \neq j \\
            \{ y_i \}_{i=0}^n \subset \mathbb{R}
        \end{cases}
    \]
    Encontrar $p \in \mathcal{P}_n$ que cumpla que:
    \[
        p(x_i) = y_i \quad \forall i = 0, 1, \ldots, n
    \]
\end{definición}

En general, el problema se puede representar como un sistema de ecuaciones, que parten de:
\[
    \begin{cases}
        y_0 = p_n(x_0) = a_0 + a_1 x_0 + a_2 x_0^2 + \ldots + a_n x_0^n \\
        y_1 = p_n(x_1) = a_0 + a_1 x_1 + a_2 x_1^2 + \ldots + a_n x_1^n \\
        \vdots \\
        y_n = p_n(x_n) = a_0 + a_1 x_n + a_2 x_n^2 + \ldots + a_n x_n^n
    \end{cases}
    \implies
\]
\[
    \underbrace{X}_{\mathcal{M}_{(n+1) \times (n+1)}} \cdot \underbrace{a}_{\mathbb{R}^{n+1}} = \underbrace{y}_{\mathbb{R}^{n+1}} \iff
    \begin{bmatrix}
        1 & x_0 & x_0^2 & \ldots & x_0^n \\
        1 & x_1 & x_1^2 & \ldots & x_1^n \\
        \vdots & \vdots & \vdots & \ddots & \vdots \\
        1 & x_n & x_n^2 & \ldots & x_n^n
    \end{bmatrix}
    \cdot
    \begin{bmatrix}
        a_0 \\
        a_1 \\
        \vdots \\
        a_n
    \end{bmatrix}
    =
    \begin{bmatrix}
        y_0 \\
        y_1 \\
        \vdots \\
        y_n
    \end{bmatrix}
\]

Y el problema ahora consistira en ver si la matriz $X$ es invertible, lo que nos deja una unica solución:
\[
    \exists ! a = X^{-1} y \implies
    \exists ! p \in \mathcal{P}_n : p = a_0 + a_1 x + a_2 x^2 + \ldots + a_n x^n
\]

\ejemplo{
    Veamos un ejemplo particular:
    \begin{center}
        \begin{tabular}{|c|ccc|}
            \hline
            $i$ & $0$ & 1 & 2 \\
            \hline
            $x$ & $0$ & $1$ & $3$ \\
            \hline
            $y$ & $-3$ & $4$ & $5$ \\
            \hline
        \end{tabular}
    \end{center}

    En este caso, el sistema de ecuaciones es el siguiente:
    \[
        \begin{cases}
            -3 = a_0 + a_1 \cdot 0 + a_2 \cdot 0^2 \\
            4 = a_0 + a_1 \cdot 1 + a_2 \cdot 1^2 \\
            5 = a_0 + a_1 \cdot 3 + a_2 \cdot 3^2
        \end{cases}
        \implies
        \begin{bmatrix}
            1 & 0 & 0 \\
            1 & 1 & 1 \\
            1 & 3 & 9
        \end{bmatrix}
        \cdot
        \begin{bmatrix}
            a_0 \\
            a_1 \\
            a_2
        \end{bmatrix}
        =
        \begin{bmatrix}
            -3 \\
            4 \\
            5
        \end{bmatrix}
    \]
}

Ahora veamos que ocurre con la matriz $X$. Para que sea invertible, es necesario que su determinante sea distinto de cero. El determinante de la matriz $X$ es el siguiente (usando la regla de Vandermonde):
\[
    \det(X) = \prod_{i \neq j} (x_i - x_j) \neq 0 \implies
    \boxed{x_i \neq x_j \iff i \neq j}
\]

\subsection{Polinomios de Lagrange}

En esta sección, vamos a ver un tipo de polinomios que nos van a resultar muy útiles para resultados teóricos.

\begin{definición}[Polinomios de Lagrange]
    Dado un conjunto de puntos $\{ x_i, y_i \}_{i=0}^n$, que nos definan un problema de interpolación polinómica, buscaremos un polinomio $p \in \mathcal{P}_n$ de la siguiente forma:
    \[
        p(x) = \sum_{i=0}^n y_i L_i(x)
    \]
    Lo que nos lleva al siguiente sistema de ecuaciones:
    \[
        \begin{cases}
            y_0 = p(x_0) = y_0 L_0(x_0) + y_1 L_1(x_0) + \ldots + y_n L_n(x_0) \\
            y_1 = p(x_1) = y_0 L_0(x_1) + y_1 L_1(x_1) + \ldots + y_n L_n(x_1) \\
            \vdots \\
            y_n = p(x_n) = y_0 L_0(x_n) + y_1 L_1(x_n) + \ldots + y_n L_n(x_n)
        \end{cases}
        \implies
        \begin{cases}
            y_0 = y_0 L_0(x_0) \\
            y_1 = y_1 L_1(x_1) \\
            \vdots \\
            y_n = y_n L_n(x_n)
        \end{cases}
    \]
    Lo que nos lleva a que:
    \[
        L_k(x_j) = \delta_{kj} = 
        \begin{cases}
            1 & \text{si } k = j \\
            0 & \text{si } k \neq j
        \end{cases}
        \implies
        L_k(x) = \prod_{j \neq k} \frac{x - x_j}{x_k - x_j}
    \]
    Donde los numeradores de la fracción se encargan de que el polinomio se anule en los puntos $x_j$ con $j \neq k$, y el denominador se encarga de que el polinomio tome el valor $1$ en el punto $x_k$ (normalizando su valor).
\end{definición}

El utilizar los polinomios de Lagrange se conoce como \textbf{Interpolación de Lagrange}.

Todo esto sera útil para interpolar funciones, creando los puntos de la forma $\{ x_i, f(x_i) \}_{i=0}^n$, y buscando un polinomio $p$ que interpole a la función $f$ en esos puntos, es decir, que cumpla que $p(x_i) = f(x_i)$ para todo $i = 0, 1, \ldots, n$.